% ============================================================================
\clearpage
\section*{v4.9.0-r56 Source-Only: Calibrating the Readout Layer}
\addcontentsline{toc}{section}{v4.9.0-r56 Source-Only: Calibrating the Readout Layer}
\label{sec:v490-r56-readout-calibration}
% ============================================================================

\paragraph{Purpose.}
The author proposed to read the structure as a core with a readout layer on top, to tune that layer until it reproduces measured values, and to learn from
the settings where they come from.  This layer records the pre-registered test H40 (note \texttt{llm/Vorab\_Notiz\_H40\_Kalibrierung\_Ausleseschicht.md},
commit \texttt{fdb27ae}, fixed before computing).  Earlier layers are not rewritten.

\noindent\textbf{Source-only status.}
This v4.9.0-r56 layer is source-only.  It does not build a PDF, does not assign a new DOI,
and does not authorize a public release.  The v4.9.0 Version DOI
\href{https://doi.org/10.5281/zenodo.23534954}{10.5281/zenodo.23534954} and the Concept DOI
\href{https://doi.org/10.5281/zenodo.19210728}{10.5281/zenodo.19210728} are unchanged by it.

\subsection*{v4.9.0-r56 the readout layer}
\label{subsec:v490-r56-layer}

For each sector \(X\in\{u,d,e,\nu\}\) a circulant coupling \(C_X=a_X+b_XP+\bar b_XP^2\) (Brannen form), a mask strength \(\varepsilon_X\) and form S of
v4.9.0-r52, \(M_X=DC_XD^{-1}\) with \(D=\exp(\varepsilon_XW_X)\).  This gives \(16\) settings (and \(m_1\)) for \(17\) measured quantities.  The admissible
settings are those of v4.9.0-r54 (CKM within \(\pm20\%\), PMNS within \(3\sigma\)).  Sixteen quantities (the four \(\varepsilon\), four ratios, and
\(|b|/|a|\), \(\theta\) per sector) were compared with fourteen candidate values from the corpus (\(1/3\), \(1/2\), \(2/3\), \(3/4\), \(1\), \(1/\sqrt3\), \(1/\sqrt2\),
\(\sqrt3\), \(2/9\), \(\pi/6\), \(\pi/9\), \(\beta_{\rm sym}^2\), \(\beta_{\rm sym}\), \(\gamma_L\)), with a control of \(10\,000\) random sets.

\subsection*{v4.9.0-r56 results}
\label{subsec:v490-r56-results}

\paragraph{Theoremlet v4.9.0-r56.1 (numerical).}
\begin{enumerate}[leftmargin=2em]
\item[(a)] Calibrated ranges: \(\varepsilon_u\in[1.36,1.63]\), \(\varepsilon_d\in[1.42,1.47]\), \(\varepsilon_e\in[0.44,0.82]\), \(\varepsilon_\nu\in[0.27,0.55]\);
\(|b|/|a|\) in \([0.953,1.032]\) (\(u\)), \([0.994,1.076]\) (\(d\)), \([1.22,6.41]\) (\(e\)), \([0.45,1.57]\) (\(\nu\)); the phases \(\theta\) fill almost all of
\([0,\pi/3]\).
\item[(b)] The candidates fall into these ranges \(76\) times; random sets do so \(65.3\) times on average, and \(1.9\%\) of them at least \(76\) times.
\item[(c)] \(|b|/|a|\approx1\) for \(u\) and \(d\) follows from the hierarchy: for eigenvalues near \((0,0,1)\) the circulant is near \(\tfrac13(1+P+P^2)\).
\end{enumerate}

\paragraph{Proof sketch.}
Script \texttt{llm/k96\_test\_h40.py}, output \texttt{llm/k96\_out.txt}; (c) by direct computation.  \(\square\)

\paragraph{Verdict and reading.}
H40a is refuted (\(p=0.019\), the pre-registered threshold is \(0.01\)); H40b (a ratio of \(\varepsilon\) containing \(1\)) holds but counts only with H40a.
\begin{enumerate}[leftmargin=2em]
\item As far as the data fix them, the settings show no structure beyond chance.  With \(16\) settings for \(17\) quantities the data hardly fix the
settings, so no structure can show in them.  The narrow quark value \(|b|/|a|\approx1\) is a property of hierarchical spectra, not of the structure.
\item For the calibration to teach something, the layer needs fewer settings (a rule fixing the Brannen quantities, or readouts by subgroups as in
v4.9.0-r55) or sharper data (CP phases, more precise masses).
\item Readings after the computation, not part of the criterion: with \(m_0\leftrightarrow a\), \(\beta\leftrightarrow|b|\), \(\gamma\leftrightarrow\varepsilon\) in
\(\Omega\) of Theorem~\ref{thm:master-field-equation}, the quarks give \(\beta/m_0\approx1\) and \(\gamma\approx1.4\)--\(1.6\), the leptons \(\gamma\approx0.3\)--\(0.8\).
Read as a deformation of the complex structure with dilatation \(K=\exp(\varepsilon\,\operatorname{span}W)\), the Beltrami modulus is
\(|\mu|\approx0.88\)--\(0.93\) (\(u\)), \(0.89\)--\(0.90\) (\(d\)), \(0.58\)--\(0.84\) (\(e\)), \(0.38\)--\(0.68\) (\(\nu\)).
\end{enumerate}

\subsection*{v4.9.0-r56 status in the three frames}
\label{subsec:v490-r56-frames}

\begin{center}
\small
\begin{tabular}{@{}>{\raggedright\arraybackslash}p{0.30\textwidth}>{\raggedright\arraybackslash}p{0.62\textwidth}@{}}
\toprule
result & M / T / P\\
\midrule
r56.1(a)--(c) & settings of mass matrices in M, the same in T as in v4.9.0-r52; P not evaluated\\
\bottomrule
\end{tabular}
\end{center}

\subsection*{v4.9.0-r56 register entries}
\label{subsec:v490-r56-register}

\begin{center}
\small
\begin{tabular}{@{}>{\raggedright\arraybackslash}p{0.34\textwidth}>{\raggedright\arraybackslash}p{0.58\textwidth}@{}}
\toprule
statement & status after v4.9.0-r56\\
\midrule
calibrated readout layer (form S, four \(\varepsilon\)) & settings recorded; no structure beyond chance (H40, \(p=0.019\))\\
\(|b|/|a|\approx1\) for quarks & consequence of the hierarchy, not of the structure\\
path ``tune the layer, then explain the settings'' & needs fewer settings or sharper data\\
\bottomrule
\end{tabular}
\end{center}

\subsection*{v4.9.0-r56 ledger and guard update}
\label{subsec:v490-r56-ledger-guard}

\[
R_{\rm cal}^{v490\text{-}r56}=(0,1,0)
\]
for (no structure in the settings; calibration recorded; compression missing).

The layer spends no guard:
\[
R_{\rm glob}^{(v490\text{-}r56,+)}=(0,1,1,1,1),
\qquad
\Gamma_{\rm guard}^{v490\text{-}r56}=(1,1,1,1,1),
\qquad
\beta_4=1/256\ \text{unspent}.
\]
It does not imply exact finite \(C5\) zero closure, exact all-readout terminality,
\(C6/T2\), a global photon-clock/RH theorem, or global all-readout closure.
